\documentclass[10pt,a4paper]{amsart}
\usepackage[margin=19mm]{geometry}
\usepackage{amsmath,amssymb,mathtools}
\usepackage[scr=rsfs,cal=euler]{mathalfa}
\usepackage{xfrac}
\usepackage[unicode,hidelinks,bookmarks=false]{hyperref}
\newcommand{\ie}{i.e.\ }
\newcommand{\cf}{cf.\ }
\newcommand{\resp}{respectively\ }
\newcommand{\Subsection}{\subsection}
\providecommand{\nobreakdash}{\nobreak\hbox{-}}
\setlength{\emergencystretch}{2em}
\multlinegap0pt
\usepackage{bm}
\usepackage{bbm}
\usepackage{dsfont}
\usepackage{xspace}
\usepackage{cancel}
\usepackage{mathtools}
\usepackage{amsthm}
\makeatletter
\@ifundefined{theorem}{\newtheorem{theorem}{Theorem}}{}
\@ifundefined{Lemma}{\newtheorem{Lemma}[theorem]{Lemma}}{}
\makeatother
\title[The Star Product of Uniformly Random Codes]{The Star Product of Uniformly Random Codes}
\author{Johan Vester Dinesen, Ragnar Freij-Hollanti, Camilla Hollanti, Benjamin Jany, Alberto Ravagnani}
\date{Source: arXiv:2511.17236v2 (CC BY 4.0)}
\begin{document}
\maketitle

\noindent\emph{Source and licence.} The Star Product of Uniformly Random Codes. Version arXiv:2511.17236v2 (CC BY 4.0), distributed under the Creative Commons Attribution 4.0 International licence, as recorded in the arXiv licence metadata for this exact version and shown on the abstract page https://arxiv.org/abs/2511.17236v2. Full rights notice: licenses/arxiv-2511.17236-CC-BY-4.0.txt.\par\smallskip

\noindent\textit{Editorial scope.} Counted unit: the Lemma whose source label is \texttt{lem:weighteddensity} in the article (source0.tex lines 668--673, source label \texttt{lem:weighteddensity}), delivered together with the article's own complete original proof of it (source0.tex lines 675--720, one \texttt{proof} environment of 46 lines). No mathematics is written, repaired or re-derived: statement and proof are the article's own text line for line. Required context retained uncounted: lem:density (lines 640--647). Every definition, display and label of those lines is retained unchanged. Omitted from this package and not redistributed: 1--639, 648--667, 674--674, 721--1296. Nothing inside a retained excerpt is omitted or altered: every retained body is delivered exactly as the article's lines are written, so no omission receipt of any kind accompanies this package. Citations: the retained excerpts carry no citation command at all, so no bibliography entry of the article is transcribed and no reference is redistributed. Reference adaptation: none was needed and none was made; every reference inside the delivered unit resolves inside it, so no unresolved reference or citation is produced. Printed slips kept literal: no printed word, symbol, hypothesis or number of the counted statement or of its proof was altered. Layout and class: the article's own class is \texttt{article} with packages including appendix, inputenc, amsmath, amssymb, mathrsfs, mathtools, enumerate, enumitem, amsthm, booktabs, xcolor, url, hyperref, multirow; the wrapper is the lane's pinned \texttt{amsart} editorial wrapper at 10pt on A4, which restates the theorem-like environments the retained text needs and adds the article's own macro definitions that the retained text needs (none), with extra wrapper packages (bm, bbm, dsfont, xspace, cancel, mathtools). The delivered numbering is the unit's own. Assets: no retained span contains a figure environment, a graphics-inclusion command, a PDF-page inclusion, a table or a TikZ picture; no image file is copied into this package, the package declares no asset dependency and no image file is redistributed with it. Licence: version arXiv:2511.17236v2 of this article is distributed under the Creative Commons Attribution 4.0 International licence, as recorded in the arXiv licence metadata for this exact version and shown on the abstract page https://arxiv.org/abs/2511.17236v2. A full rights notice is delivered at licenses/arxiv-2511.17236-CC-BY-4.0.txt. Characters: every retained excerpt is pure ASCII, so the delivered engine, XeLaTeX with its default Latin Modern text font, reports no missing character.
\begin{Lemma}\label{lem:density}
Let $k_1,k_2\colon \mathbb{N}\rightarrow \mathbb{N}$ be strictly monotone increasing functions such that $k_1(t)\leq k_2(t)$ for all $t\in \mathbb{N}$.
Fix $\alpha\in(0,1)$. 
Then 
\begin{align*}
    \lim_{t\rightarrow \infty}\frac{|S^-(\alpha)|}{|S|} = 0 \quad \text{and} \quad \lim_{t\rightarrow \infty}\frac{|S^+(\alpha)|}{|S|} = 1.
\end{align*}
\end{Lemma}

\begin{Lemma}
Let $k_1,k_2\colon \mathbb{N}\rightarrow \mathbb{N}$ be strictly monotone increasing functions such that $k_1(t)\leq k_2(t)$ for all $t\in \mathbb{N}$.
Then for any $\alpha \in [0,1)$ we have
    $$\lim_{t \rightarrow \infty} \frac{\sum_{r=\lfloor\alpha k_1(t)\rfloor}^{k_1(t)} \sum_{\ell=0}^{k_2(t)-k_1(t)} \binom{k_2(t)-k_1(t)}{\ell}S_{r,\ell}q^{\ell}}{\sum_{r=\lfloor \alpha k_1(t) \rfloor}^{k_1(t)} \sum_{\ell=0}^{k_2(t)-k_1(t)} \binom{k_2(t)-k_1(t)}{\ell}S_{r,\ell}} = \lim_{t\rightarrow \infty} \exp \left( \frac{(q-1)k_2(t)}{q^{k_1(t)}+1}\right).
$$ \label{lem:weighteddensity}
\end{Lemma}

\begin{proof}
We begin by proving the claim for $\alpha = 0$, and then show that this implies it for any $0 \le \alpha <1$. Note that $\sum_{r=0}^{k_1(t)} S_{r,\ell} = q^{k_1(t)k_2(t) - k_1(t) - k_1(t)\ell}$. This allows us to rewrite the LHS of the statement as
    \begin{align*}
       \frac{q^{k_1(t) k_2(t) - k_1(t)}\sum_{\ell = 0}^{k_2(t)-k_1(t)} \binom{k_2(t)-k_1(t)}{\ell}q^{-\ell(k_1(t)-1)}}{q^{k_1(t)k_2(t)-k_1(t)}\sum_{\ell=0}^{k_2(t)-k_1(t)}\binom{k_2(t)-k_1(t)}{\ell}q^{-\ell k_1(t)}} &= \frac{(q^{-(k_1(t)-1)} +1)^{k_2(t)-k_1(t)}}{(q^{-k_1(t)} + 1)^{k_2(t)-k_1(t)}}\\
                &= \left(\frac{q^{k_1(t)} +q}{q^{k_1(t)}+1}\right)^{k_2(t)-k_1(t)}\\
                &= \left(1 + \frac{q-1}{q^{k_1(t)}+1}\right)^{k_2(t)-k_1(t)}, 
    \end{align*}
where the first equality follows from the Binomial Theorem. 


We distinguish two cases according to the asymptotic behavior of $k_2(t)-k_1(t)$, which will be exhaustive by the assumed properties of $k_1$ and $k_2$. For the first case suppose $\lim_{t \rightarrow \infty} k_2(t) - k_1(t) =\infty$. It is well known that for functions $h,p\colon \mathbb{R}\rightarrow \mathbb{R}$ such that $\lim_{t \rightarrow \infty} h(t) = 0$ and $\lim_{t \rightarrow \infty} p(t) = \infty$, we have $$\lim_{t \rightarrow \infty}(1+h(t))^{p(t)} = \lim_{t \rightarrow \infty} \exp({h(t)p(t)}).$$
Hence, letting $h(t) \coloneqq  \frac{q-1}{q^{k_1(t)}+1}$ and $p(t) \coloneqq k_2(t)-k_1(t)$, 
we get the desired identity for $\alpha = 0$ and $\lim_{t \rightarrow \infty} k_2(t) - k_1(t) =\infty$.


For the second case, suppose $\lim_{t \rightarrow \infty} k_2(t) - k_1(t) = C < \infty.$ Since $\lim_{t \rightarrow \infty} k_1(t) = \infty$ the result follows as
\begin{align*}
    1= \lim_{t \rightarrow \infty} \left(1 + \frac{q-1}{q^{k_1(t)}+1}\right)^{k_2(t)-k_1(t)} &=  \lim_{t \rightarrow \infty} \exp\left({\frac{q-1}{q^{k_1(t)}+1}(k_2(t)-k_1(t))}\right). \\
    &=  \lim_{t \rightarrow \infty} \exp\left({\frac{(q-1)k_2(t)}{q^{k_1(t)}+1}}\right).
\end{align*}

Suppose now $0<\alpha <1$ and let $$N = \sum_{r=0}^{k_1(t)} \sum_{\ell=0}^{k_2(t)-k_1(t)} \binom{k_2(t)-k_1(t)}{\ell}S_{r,\ell}q^{\ell},$$ and $$D =\sum_{r=0}^{k_1(t)} \sum_{\ell=0}^{k_2(t)-k_1(t)} \binom{k_2(t)-k_1(t)}{\ell}S_{r,\ell} = |S|.$$ Additionally, partition the sums of $D$ and $N$ as follows:
\begin{align*}
    N' &= \sum_{r=\lfloor\alpha k_1(t)\rfloor}^{k_1(t)} \sum_{\ell=0}^{k_2(t)-k_1(t)} \binom{k_2(t)-k_1(t)}{\ell}S_{r,\ell}q^{\ell},\\
    \Delta N &= \sum_{r=0}^{\lfloor \alpha k_1(t) \rfloor -1} \sum_{\ell=0}^{k_2(t)-k_1(t)} \binom{k_2(t)-k_1(t)}{\ell}S_{r,\ell}q^{\ell}, \\
    D' &= \sum_{r=\lfloor \alpha k_1(t) \rfloor}^{k_1(t)} \sum_{\ell=0}^{k_2(t)-k_1(t)} \binom{k_2(t)-k_1(t)}{\ell}S_{r,\ell},\\
    \Delta D &= \sum_{r=0}^{\lfloor \alpha k_1(t) \rfloor -1} \sum_{\ell=0}^{k_2(t)-k_1(t)} \binom{k_2(t)-k_1(t)}{\ell}S_{r,\ell}. \\
\end{align*}
Note that when $r=0$, we have $\sum_{\ell=0}^{k_2(t)-k_1(t)} \binom{k_2(t)-k_1(t)}{\ell}S_{0,\ell} =1$. Hence, when considering  $\frac{\Delta D -1}{D}$, we can apply Lemma \ref{lem:density} and its proof to obtain  $\lim_{t \rightarrow \infty} \frac{\Delta D}{D} = \lim_{t \rightarrow \infty} \frac{\Delta D - 1}{D} =  0$ as 
\[
    \frac{\Delta D -1}{D} \leq q^{(\alpha-1) k_1(t) k_2(t) + (\alpha-\alpha^2) k_1(t)^2 + \alpha k_1(t) + \log_q(\alpha k_1(t))}.
\]
Furthermore, as $1\leq q^\ell \leq q^{k_2(t)-k_1(t)}$ for all $\ell \leq k_2(t)-k_1(t)$, we have
\begin{align*}
    \frac{\Delta N - 1}{N} &\leq q^{k_2(t)-k_1(t)} \frac{\Delta D -1}{D} \\
    &\leq q^{(\alpha-1) k_1(t) k_2(t) + (\alpha-\alpha^2) k_1(t)^2 + k_2(t) + (\alpha -1)k_1(t) + \log_q(\alpha k_1(t))}.
\end{align*}
This allows us to conclude that $\lim_{t\rightarrow \infty} \frac{\Delta N}{N} = \lim_{t\rightarrow \infty} \frac{\Delta N -1 }{N} = 0$. The result then follows as
\begin{align*}
   \lim_{t\rightarrow \infty} \frac{N'}{D'} &= \lim_{t\rightarrow \infty}\frac{N-\Delta N}{D- \Delta D}\\
   &= \lim_{t\rightarrow \infty} \frac{N}{D} \left(\frac{1- \frac{\Delta N}{N}}{1- \frac{\Delta D}{D}} \right)  \\
   &= \lim_{t\rightarrow \infty} \frac{N}{D}\\ 
     &= \lim_{t\rightarrow \infty} \exp \left( \frac{(q-1)k_2(t)}{q^{k_1(t)}+1}\right)
  \end{align*}
where the last equality follows from the case $\alpha = 0$. 
\end{proof}

\end{document}
